\section{Graph-Theoretic Interpretation}

There are three useful graph-theoretic layers.

\paragraph{The Cayley graph layer.}
Let
\[
  \XN = \mathrm{Cay}(S_N,T_N).
\]
This is the undirected $\binom{N}{2}$-regular Cayley graph generated by all
transpositions.  A word $w=(t_1,\ldots,t_m)$ determines a walk
\[
  x_0=\id,\qquad x_j=t_1t_2\cdots t_j,\qquad 0\le j\le m,
\]
in $\XN$.  The identity-word condition is exactly $x_m=\id$.  Thus
$\Omega(N,m)$ is the set of length-$m$ closed walks in $\XN$ rooted at
$\id$, and the target law is the uniform closed-walk bridge measure.  The
intractable exact bridge sampler is the usual endpoint-conditioned random walk
on $\XN$.

\paragraph{The complete-graph edge-sequence layer.}
Since transpositions are edges of the complete graph $\KN$ on vertex set
$[N]$, a word is also an ordered edge sequence
\[
  e_1,\ldots,e_m\in E(\KN).
\]
Evaluating the word is the token-swapping process that swaps the tokens at the
endpoints of $e_j$ at time $j$.  Hence $\Omega(N,m)$ is the set of length-$m$
ordered edge sequences in $\KN$ whose token-swapping schedule returns every
token to its starting vertex.

For a word $w$, let $G(w)$ be the support multigraph on $[N]$ whose edge
multiset is $\{e_1,\ldots,e_m\}$, with the time order retained separately.
The identity condition is not equivalent to a simple parity condition on
$G(w)$: even edge multiplicities or even degrees are not sufficient because
the swaps do not commute in general.
However, $G(w)$ still records important coarse geometry.  Its connected
components, together with the number of time steps supported in each
component, give the weighted support-component data used throughout the
experiments.  Disjoint support components commute, so after sorting factors
with disjoint supports, the identity condition decomposes over these
components.

\paragraph{The rewrite-graph layer.}
Each sampler defines a graph whose vertices are the closed schedules
$\Omega(N,m)$ and whose edges are legal local rewrites.  Pure Hurwitz
sampling is random walk on the Hurwitz rewrite graph.  Same-product
block heat-bath is random walk on a larger local path-replacement graph:
choose a contiguous time interval, keep the associated permutation of that
window fixed, and replace the interval by another ordered edge sequence with
the same endpoint in $\XN$.

In the complete-graph language, Hurwitz moves have a simple local form.  If
two adjacent edges are disjoint, the move just swaps their order.  If two
adjacent edges share a vertex, for example
\[
  (ij),(jk),
\]
then a right Hurwitz move replaces them by
\[
  (jk),(ik).
\]
Thus Hurwitz moves are generated by disjoint-edge commutations and local
triangle moves in $\KN$.  A square pair $(t,t)$ is fixed by Hurwitz moves,
whereas block heat-bath can perform square refreshes $(t,t)\mapsto (s,s)$
inside a length-$4$ same-product context.

With this vocabulary, the current state of the project is:
\begin{itemize}
\item pure Hurwitz rewrite graphs are disconnected on $\Omega(N,m)$ in small
  exact cases;
\item the conjectural input \texttt{H1\_id} says that, on one connected
  support component, every feasible identity schedule is connected by
  triangle and disjoint-commutation moves;
\item the feasible identity-component lemma says that a connected component
  on $v$ vertices can support an identity schedule of length $k$ exactly when
  $k$ is even and $k\ge 2(v-1)$, with doubled trees as the minimal examples;
\item assuming \texttt{H1\_id}, the $\{3,4,5\}$ same-product rewrite graph is
  connected for the nonexceptional even-length regimes;
\item the polynomial-mixing problem is an expansion or spectral-gap problem
  for this local rewrite graph.
\end{itemize}

The projection/decomposition proof route is therefore a graph-expansion
program: prove expansion of the projection to weighted support-component
macrostates and prove expansion inside each fiber.  The large-block comparison
route compares the local path-replacement graph to a mesoscopic
path-replacement graph.  The negative canonical-routing result is also
graph-theoretic: deterministic paths through a fixed doubled-tree or
canonical-square normal form create superpolynomial congestion because such
normal forms have exponentially small stationary mass.
